\begin{abstract}
% Length: ~5%
% Focus: Write this last. Keep it punchy and direct.

% Must contain:
% 1. Background: Why the industry relies on LLM judges for RAG.
% 2. Pain point: Long text and retrieved noise compromise the judge.
% 3. Our focus: We investigate the micro-level cognitive mechanisms of the judge, going beyond macro-level dataset lists.
% 4. Conclusion: The necessary shift from static judges to active agentic frameworks.
Retrieval Augmented Generation systems increasingly rely on Large Language Models as automated evaluators due to the unscalable nature of human annotation. However the reliability of these judges under the unique stressors of retrieval augmented environments remains under explored. This survey investigates the micro level cognitive vulnerabilities of evaluators by dissecting the internal conflict between truthfulness and faithfulness. We analyze how environmental catalysts such as extreme context lengths and retrieved noise act as specific triggers for inherent biases. Beyond documenting macro level performance drops we explore the internal mechanistic conflicts that occur during knowledge reconciliation. Furthermore we track the evolutionary trajectory of mitigation strategies moving from static interventions to multi model debates. Ultimately we demonstrate the necessity of a paradigm shift toward proactive tool augmented agentic frameworks capable of external verification. This analysis provides a robust foundation for researchers seeking to design bias resistant assessment architectures.
\end{abstract}